\documentclass[11pt]{article}

\usepackage[margin=1in]{geometry}
\usepackage{amsmath,amssymb}
\usepackage{hyperref}

\title{qanneal: A CPU-First Simulated Quantum Annealing Toolkit}
\author{qanneal contributors}
\date{\today}

\begin{document}
\maketitle

\begin{abstract}
This report describes the mathematical foundations and software design of \texttt{qanneal}, a CPU-first simulated quantum annealing (SQA) toolkit. We detail the Ising and QUBO formulations, the Suzuki--Trotter mapping for transverse-field models, and the Monte Carlo algorithms used in classical SA and SQA. We also document the code architecture, trace observability, and practical complexity considerations.
\end{abstract}

\section{Introduction}
Combinatorial optimization problems can often be written as minimizing a quadratic energy function over binary variables. Simulated annealing (SA) provides a classical stochastic method for approximately solving these problems. Simulated quantum annealing (SQA) extends this approach by emulating quantum fluctuations through a path-integral representation of a transverse-field Ising model. \texttt{qanneal} provides SA and SQA with reproducible schedules, sweep-level tracing, and a C++ core with Python bindings.

\section{Problem Formulation}

\subsection{Ising model}
Given spins $s_i \in \{-1, +1\}$, the Ising energy is
\begin{equation}
E(\mathbf{s}) = \sum_i h_i s_i + \sum_{i,j} J_{ij} s_i s_j + c.
\end{equation}
The goal is to minimize $E(\mathbf{s})$.

\subsection{QUBO}
For binary variables $x_i \in \{0,1\}$, a quadratic unconstrained binary optimization (QUBO) instance is
\begin{equation}
E(\mathbf{x}) = \sum_{i,j} Q_{ij} x_i x_j.
\end{equation}
We map QUBO to Ising with $x_i = (1 + s_i)/2$, yielding an equivalent Ising model with transformed $h$, $J$, and constant shift $c$.

\section{Simulated Annealing}
SA performs a Markov chain over spin configurations. At each temperature $T$ (or inverse temperature $\beta$), it proposes single-spin flips and accepts them with Metropolis probability
\begin{equation}
P(\text{accept}) = \min\left(1, e^{-\beta \Delta E}\right).
\end{equation}
The schedule $\beta_0,\beta_1,\ldots$ defines the anneal.

\section{Simulated Quantum Annealing}

\subsection{Transverse-field Ising model}
The quantum Hamiltonian is
\begin{equation}
\hat{H} = \sum_i h_i \hat{\sigma}^z_i + \sum_{i,j} J_{ij} \hat{\sigma}^z_i \hat{\sigma}^z_j - \Gamma \sum_i \hat{\sigma}^x_i.
\end{equation}
Here $\Gamma$ is the transverse-field strength controlling quantum fluctuations.

\subsection{Suzuki--Trotter mapping}
SQA uses a path-integral representation by discretizing imaginary time into $M$ slices. This yields a classical Ising model over $M$ replicas with nearest-neighbor coupling along the slice dimension. The effective coupling is
\begin{equation}
J_\perp = \frac{1}{2} \log \left(\coth \left(\frac{\beta \Gamma}{M}\right)\right).
\end{equation}
The algorithm then performs Monte Carlo updates over the augmented state space.

\subsection{Update phases}
\texttt{qanneal} performs two update phases per schedule step:
\begin{itemize}
\item Slice updates: single-spin flips within each slice.
\item Worldline updates: collective flips of a spin across all slices.
\end{itemize}
Both are accepted with Metropolis probability using the effective classical energy.

\section{Implementation Architecture}
The C++ core defines model classes (DenseIsing, SparseIsing, QUBO), schedulers, and annealers. The \texttt{Backend} interface isolates the annealing logic from energy and delta-energy computation. Python bindings are provided via pybind11.

\subsection{State representation}
\texttt{State} stores spins as \texttt{int8\_t} values. \texttt{SQAState} stores data in a flattened buffer with layout:
\begin{equation}
\texttt{index} = (r \cdot M + t) \cdot N + i
\end{equation}
where $r$ is replica, $t$ is slice, and $i$ is spin.

\subsection{Observability}
Sweep-level observers record the full state, energy traces, and replica-wise energies. This enables reproducibility, diagnostics, and post-hoc analysis of the anneal.

\section{Complexity Considerations}
Let $N$ be the number of spins, $M$ the number of slices, $R$ the number of replicas, and $S$ the number of sweeps per step. For dense models, each sweep costs $O(N^2)$; for sparse models, the cost scales with the number of edges. SQA multiplies compute by $M \cdot R$.

\section{Experimental Protocol (Template)}
To compare algorithms:
\begin{itemize}
\item Fix a problem family and size $N$.
\item Use identical schedules and sweep budgets.
\item Run multiple seeds and report distributions of energies.
\item Use sweep-level traces to study convergence and mixing.
\end{itemize}

\section{Limitations}
\begin{itemize}
\item CPU-first; CUDA is a planned extension.
\item Dense models scale quadratically in memory and time.
\item SQA is sensitive to schedule tuning and slice count.
\end{itemize}

\section{Future Work}
\begin{itemize}
\item CUDA backend for dense and sparse kernels.
\item Cluster updates and improved worldline moves.
\item MPI-based multi-node replicas.
\item Automated schedule tuning and diagnostics.
\end{itemize}

\section{Reproducibility}
Each annealer supports explicit RNG seeds. Sweep-level traces can be saved for deterministic replay and analysis. The project ships example scripts and unit tests for baseline verification.

\end{document}
